\paragraph{Observações.}
\begin{itemize}
  \item \textbf{Com as hipóteses do artigo (iii), o FS mínimo reproduz a Eq.~6 nos nove casos:} as diferenças vão de
    $-0{,}016$ (f50, 36~mm/h: 1,034 contra 1,05) a $+0{,}021$ (f10, 90~mm/h: 1,071 contra 1,05). A Eq.~6 é um ajuste
    aos pontos do artigo, então diferenças de 0,02 estão no nível da própria Eq.~6.
  \item \textbf{A Fig.~5b é reproduzida ponto a ponto} (Tabela~\ref{tab:fig5b}): as doze diferenças ficam entre
    $-0{,}028$ e $+0{,}028$, dentro da digitalização ($\pm 0{,}03$).
  \item \textbf{O relatório de outubro ficava acima do artigo nas chuvas moderadas}, até $+0{,}17$ (f50, 36~mm/h, 12~h:
    1,57 contra 1,40) e $+0{,}14$ no FS mínimo do mesmo caso (1,18 contra 1,05, contra 1,034 hoje). O relatório
    atribuía a diferença ao SRM contra Bishop e à digitalização; com o SRM de hoje ela desaparece, então a causa
    estava no cálculo antigo. Candidatos: a chuva (hoje por metro inclinado, o relatório na projeção horizontal: na
    face de $45^\circ$ hoje entra 41\,\% a mais de água) e a malha do transiente (as duas são testadas no fermi, T4).
    Nas chuvas fortes, onde a superfície alaga e a quantidade de chuva aplicada deixa de importar, os dois concordam
    (f10 180: 0,842 e 0,86; f50 90: 0,832 e 0,85), o que é coerente com a convenção da chuva como causa principal.
\end{itemize}
